%! Author = ben
%! Date = 27.11.2023

% Preamble
\documentclass[../main.tex]{subfiles}

% Packages

% Document
\begin{document}
    \chapter{Vorgehensweisen, Materialien und Methoden}
    SaveWorld besteht aus mehreren Komponenten, die alle miteinander interagieren.
    Die beiden Hauptkomponenten sind die App (Frontend) und der Server (Backend).
    F\"ur die Verwirklichung der App wird das Framework Ionic mit einer angepassten Design Bibliothek verwendet.
    Ionic ist ein Framework für sogenante Hybrid-Apps, welche auf allen gängigen Plattformen laufen.
    Das Backend wurde mit dem Express Framework für Node.js entwickelt.
    Die Kommunikation zwischen Backend und Frontend verläuft zum Großteil über eine REST-API\footnote{
        REST steht für \dq Representational State Transfer\dq\ und ist ein Paradigma der Softwarearchitektur,
        welches das Verhalten des World Wide Webs abstrahiert, damit es den Anforderungen des modernen Webs besser genügt.
    }.
    Damit die Verwendung dieser REST-API in der App einfacher und einheitlicher ist,
    wurde ein Software Development Kit (SDK) für die Programmiersprache JavaScript entwickelt.
\end{document}
